% Course Lecture Template Sample
% 课程讲座模板样板
% 适用于：大学课程、企业培训、在线教学、技术分享

\documentclass[aspectratio=43,12pt]{beamer}

% ========== 基础包 ==========
\usepackage[fontset=fandol,space=true,scheme=plain]{ctex}
\usetheme{Madrid}
\usecolortheme{crane}

% ========== 额外包 ==========
\usepackage{graphicx}
\usepackage{amsmath,amssymb}
\usepackage{tikz}
\usetikzlibrary{shapes,arrows,positioning}
\usepackage{listings}
\usepackage{xcolor}

% ========== 代码高亮设置 ==========
\lstset{
    language=Python,
    basicstyle=\ttfamily\small,
    keywordstyle=\color{blue},
    commentstyle=\color{green!60!black},
    stringstyle=\color{red},
    numbers=left,
    numberstyle=\tiny\color{gray},
    frame=single,
    breaklines=true
}

% ========== 文档信息 ==========
\title{\textbf{Python 编程基础}}
\subtitle{第 1 课：变量与数据类型}
\author{讲师姓名}
\institute{XX大学计算机学院}
\date{\today}

% ========== 开始文档 ==========
\begin{document}

% ---------- 封面 ----------
\begin{frame}
    \titlepage
\end{frame}

% ---------- 课程概述 ----------
\begin{frame}{课程概述}
    本课程是 Python 编程的入门课程，共 12 课时。
    
    \vspace{0.5em}
    \textbf{主要内容：}
    \begin{itemize}
        \item Python 基础语法
        \item 数据类型与结构
        \item 流程控制语句
        \item 函数与模块
        \item 文件操作与异常处理
        \item 面向对象编程基础
    \end{itemize}
    
    \vspace{0.5em}
    \textbf{考核方式：}平时作业 40\% + 期末项目 60\%
\end{frame}

% ---------- 学习目标 ----------
\begin{frame}{本节课学习目标}
    学完本节课后，你将能够：
    \begin{enumerate}
        \item 理解什么是变量以及变量的作用
        \item 掌握 Python 的基本数据类型（整数、浮点数、字符串、布尔值）
        \item 学会正确的变量命名规则
        \item 完成简单的输入输出操作
    \end{enumerate}
    
    \vspace{1em}
    \begin{block}{预计时长}
        45 分钟讲解 + 15 分钟练习
    \end{block}
\end{frame}

% ---------- 知识点：什么是变量 ----------
\begin{frame}{什么是变量}
    \begin{block}{定义}
        变量是用于存储数据的容器。可以将变量想象成一个贴有标签的盒子，
        盒子里装着数据，标签就是变量名。
    \end{block}
    
    \vspace{1em}
    \textbf{类比理解：}
    \begin{itemize}
        \item 变量名 = 标签（如"年龄"、"分数"）
        \item 变量值 = 盒子里的内容（如 20、95.5）
        \item 数据类型 = 盒子的类型（装数字的、装文字的）
    \end{itemize}
\end{frame}

% ---------- 知识点：数据类型 ----------
\begin{frame}{Python 的基本数据类型}
    \begin{columns}[T]
        \begin{column}{0.48\textwidth}
            \textbf{数值类型}
            \begin{itemize}
                \item \texttt{int} - 整数\\
                例：\texttt{age = 20}
                \item \texttt{float} - 浮点数\\
                例：\texttt{score = 95.5}
                \item \texttt{complex} - 复数\\
                例：\texttt{x = 3 + 4j}
            \end{itemize}
        \end{column}
        \begin{column}{0.48\textwidth}
            \textbf{其他类型}
            \begin{itemize}
                \item \texttt{str} - 字符串\\
                例：\texttt{name = "张三"}
                \item \texttt{bool} - 布尔值\\
                例：\texttt{passed = True}
                \item \texttt{None} - 空值\\
                例：\texttt{result = None}
            \end{itemize}
        \end{column}
    \end{columns}
\end{frame}

% ---------- 示例演示 ----------
\begin{frame}[fragile]{示例：变量定义与使用}
    \begin{lstlisting}
# 定义变量
name = "张三"       # 字符串
age = 20           # 整数
score = 95.5       # 浮点数
is_student = True  # 布尔值

# 查看变量类型
print(type(name))   # <class 'str'>
print(type(age))    # <class 'int'>
print(type(score))  # <class 'float'>

# 使用变量
print(name, "的年龄是", age)
print("成绩:", score)
    \end{lstlisting}
\end{frame}

% ---------- 命名规则 ----------
\begin{frame}{变量命名规则}
    \begin{columns}[T]
        \begin{column}{0.55\textwidth}
            \textbf{必须遵守的规则：}
            \begin{enumerate}
                \item 只能包含字母、数字、下划线
                \item 不能以数字开头
                \item 不能是 Python 关键字
                \item 区分大小写
            \end{enumerate}
            
            \vspace{0.5em}
            \textbf{建议遵守的规范：}
            \begin{itemize}
                \item 使用有意义的名称
                \item 小写字母加下划线
            \end{itemize}
        \end{column}
        \begin{column}{0.43\textwidth}
            \begin{exampleblock}{命名示例}
                \textcolor{green!60!black}{\cmark} \texttt{user\_name}\\
                \textcolor{green!60!black}{\cmark} \texttt{score2024}\\
                \textcolor{green!60!black}{\cmark} \texttt{\_temp}\\
                \vspace{0.3em}
                \textcolor{red}{\xmark} \texttt{2nd\_name}\\
                \textcolor{red}{\xmark} \texttt{my-name}\\
                \textcolor{red}{\xmark} \texttt{class}
            \end{exampleblock}
        \end{column}
    \end{columns}
\end{frame}

% ---------- 课堂练习 ----------
\begin{frame}{课堂练习}
    \begin{exampleblock}{练习题}
        请写出以下代码的输出结果：
        \begin{lstlisting}
x = 10
y = 20
print(x + y)
print(x * y)
        \end{lstlisting}
    \end{exampleblock}
    
    \pause
    \vspace{1em}
    \textbf{答案：}
    \begin{itemize}
        \item \texttt{30} （10 + 20 的结果）
        \item \texttt{200} （10 × 20 的结果）
    \end{itemize}
\end{frame}

% ---------- 小结 ----------
\begin{frame}{本节课小结}
    \begin{columns}[T]
        \begin{column}{0.31\textwidth}
            \textbf{变量}
            \begin{itemize}
                \item 存储数据的容器
                \item 变量名 + 变量值
                \item 命名有规则
            \end{itemize}
        \end{column}
        \begin{column}{0.31\textwidth}
            \textbf{数据类型}
            \begin{itemize}
                \item int - 整数
                \item float - 小数
                \item str - 字符串
                \item bool - 布尔值
            \end{itemize}
        \end{column}
        \begin{column}{0.31\textwidth}
            \textbf{输入输出}
            \begin{itemize}
                \item \texttt{print()} 输出
                \item \texttt{input()} 输入
                \item \texttt{type()} 查类型
            \end{itemize}
        \end{column}
    \end{columns}
\end{frame}

% ---------- 作业布置 ----------
\begin{frame}{课后作业}
    \textbf{必做题：}
    \begin{enumerate}
        \item 定义三个变量分别存储你的姓名、年龄和爱好，并打印出来
        \item 编写程序计算圆的面积（半径 5）
        \item 指出以下变量名的正误：\texttt{my-var}、\texttt{\_name}、\texttt{2nd}
    \end{enumerate}
    
    \vspace{1em}
    \textbf{选做题：}
    \begin{itemize}
        \item 探索 Python 的其他数据类型（list、dict、tuple）
        \item 尝试使用 \texttt{input()} 实现交互式程序
    \end{itemize}
    
    \vspace{0.5em}
    \textbf{提交截止时间：}下周上课前
\end{frame}

% ---------- 结束页 ----------
\begin{frame}
    \centering
    \Huge{\textbf{谢谢！}}
    
    \vspace{2em}
    \large{Q \& A}
    
    \vspace{1em}
    \normalsize
    如有问题，请联系：\\
    \texttt{teacher@university.edu.cn}
\end{frame}

\end{document}
